\medskip

On considère la fonction $f$ définie sur $]0 ~;~+ \infty[$ par 

\[f(x) = x^2 - x \ln (x).\]

On admet que $f$ est deux fois dérivable sur $]0 ~;~+ \infty[$.

On note $f'$ la fonction dérivée de la fonction $f$ et $f''$ la fonction dérivée de la fonction $f'$.

\medskip

\textbf{Partie A :} Étude de la fonction $f$

\medskip

\begin{enumerate}
\item Déterminer les limites de la fonction $f$ en $0$ et en $+\infty$.
\item Pour tout réel $x$ strictement positif, calculer $f'(x)$.
\item Montrer que pour tout réel $x$ strictement positif:

\[f''(x) = \dfrac{2x - 1}{x}.\]

\item Étudier les variations de la fonction $f'$ sur $]0 ~;~+ \infty[$, puis dresser le tableau des variations de la fonction $f'$ sur $]0 ~;~+ \infty[$.

On veillera à faire apparaître la valeur exacte de l'extremum de la fonction $f'$ sur $]0 ~;~+ \infty[$.

Les limites de la fonction $f'$ aux bornes de l'intervalle de définition ne sont pas attendues.
\item Montrer que la fonction $f$ est strictement croissante sur $]0 ~;~+ \infty[$.
\end{enumerate}

\medskip

\textbf{Partie B :} Étude d'une fonction auxiliaire pour la résolution de l'équation $f(x) = x $

\medskip

On considère dans cette partie la fonction $g$ définie sur $]0 ~;~+ \infty[$ par 

\[g(x) = x - \ln (x).\]

On admet que la fonction $g$ est dérivable sur $]0 ~;~+ \infty[$, on note $g'$ sa dérivée.

\medskip

\begin{enumerate}
\item Pour tout réel strictement positif, calculer $g'(x)$, puis dresser le tableau des variations de la fonction $g$.

Les limites de la fonction $g$ aux bornes de l'intervalle de définition ne sont pas attendues.
\item On admet que 1 est l'unique solution de l'équation $g(x) = 1$.

Résoudre, sur l'intervalle $]0 ~;~+ \infty[$, l'équation $f(x) = x$.
\end{enumerate}

\medskip

\textbf{Partie C :} Étude d'une suite récurrente

On considère la suite $\left(u_n\right)$ définie par $u_0 = \dfrac12$ et pour tout entier naturel $n$,

\[u_{n+1} = f\left(u_n\right) = u_n^2  - u_n \ln \left(u_n\right).\]

\medskip

\begin{enumerate}
\item Montrer par récurrence que pour tout entier naturel $n$:
\[\dfrac12  \leqslant u_n \leqslant u_{n+1} \leqslant 1.\]

\item Justifier que la suite $\left(u_n\right)$ converge.

On appelle $\ell$ la limite de la suite $\left(u_n\right)$ et on admet que $\ell$ vérifie l'égalité $f(\ell) = \ell$.
\item Déterminer la valeur de $\ell$.
\end{enumerate}

\bigskip
